
\exo 

déterminer l'équation réduite de la droite $(AB)$ avec $A(2 ; 4)$ et $B(4 ; -3)$

\exo

Déterminer une équation cartésienne de la droite $(AB)$ avec $A(3 ; -5)$ et $B(7 ; -1)$ en utilisant un déterminant de deux vecteurs.

\exo 


Donner un vecteur directeur des droites suivantes :

\begin{enumerate}
\item $D_1 : 3x + 5y -7 = 0$ \dotfill
\item $D_2 : 4x -y +9 = 0$ \dotfill
\item $ D_3 : x = 6$ \dotfill
\item $D_4 : y = -5x +8$ \dotfill
\item $D_5 :  x + y -3=0$ \dotfill
\end{enumerate}

\exo 

Les droites $\mathscr{D} : 3x -6x+ 8=0$ et $\mathscr{D'} : -2x + 4x -6=0$ sont-elles parallèles ?


\exo

Déterminer une équation cartésienne de la droite passant par $A(1;2)$ et parallèle à la droite $(BC)$ où $B(4; -2)$ et $C(7; -1)$.

\exo

Résoudre, à l'aide de la méthode par substitution, le système : $\left\lbrace \begin{array}{c@{~=~}c} 
x+4y&-7\\
2x+5y&-8
\end{array}\right. $

\exo

Résoudre, à l'aide de la méthode par combinaisons linéaires, le système : $\left\lbrace \begin{array}{c@{~=~}c} 
5x+4y&7\\
2x-3y&12
\end{array}\right. $

\exo

A la fin de la saison de la chasse, l'on demande à Nemrod combien il a tué de lièvres et de faisans.

\og 19 têtes et 54 pattes \fg{} 

\begin{enumerate}
\item Traduire le texte sous forme d'équations.

{\it Indication : Un lièvre a une seule tête et quatre pattes tandis qu'un faisan a une tête et deux pattes.}

\item Résoudre le système obtenu et conclure.
\end{enumerate}

\exo \qquad \og L’âne et le mulet \fg{} (extraits de textes de Léonhard Euler 1707-1783)


Un mulet et un âne portent des charges de quelques quintaux. L'âne se plaint de la sienne et dit au mulet : \og il ne me manque que de porter encore un quintal de ta charge pour être alors chargé du double que tu ne le serais. \fg{}

Le mulet répond : \og Oui, mais si tu me donnais un quintal de la tienne, je serais trois fois plus chargé que tu ne le serais.  \fg{}

Combien de quintaux portent-ils chacun ?
